% !TeX program = lualatex
\documentclass[12pt,aspectratio=169]{beamer}
\usepackage[utf8]{inputenc}
\usepackage[T1]{fontenc}
\usepackage{amsmath,amsfonts,amssymb}
\usepackage{graphicx}
\usepackage{xcolor}
\usepackage{hyperref}
\usepackage{anyfontsize}
\usepackage{lmodern}
\usepackage{longtable}
\usepackage{array}
\usepackage{booktabs}

% ====== EXTREME FONT REDUCTION FOR MAX CONTENT ======
\renewcommand{\tiny}{\fontsize{3.5}{4.5}\selectfont}
\renewcommand{\scriptsize}{\fontsize{4.5}{5.5}\selectfont}
\renewcommand{\footnotesize}{\fontsize{5.5}{6.5}\selectfont}
\renewcommand{\small}{\fontsize{6.5}{7.5}\selectfont}
\renewcommand{\normalsize}{\fontsize{7.5}{8.5}\selectfont}
\renewcommand{\large}{\fontsize{8.5}{9.5}\selectfont}
\renewcommand{\Large}{\fontsize{9.5}{10.5}\selectfont}
\renewcommand{\LARGE}{\fontsize{10.5}{12}\selectfont}
\renewcommand{\huge}{\fontsize{11.5}{13.5}\selectfont}
\renewcommand{\Huge}{\fontsize{13}{16}\selectfont}

% ====== COLORS ======
\definecolor{redcorrect}{RGB}{200,30,30}
\definecolor{bluecorrect}{RGB}{30,80,200}
\definecolor{greencheck}{RGB}{0,150,50}
\definecolor{lightgray}{RGB}{245,245,245}
\definecolor{darkgray}{RGB}{40,40,40}
\definecolor{softyellow}{RGB}{255,248,200}
\definecolor{deepblue}{RGB}{20,60,150}
\definecolor{darkred}{RGB}{150,20,20}
\definecolor{orangewarning}{RGB}{255,140,0}
\definecolor{correctgreen}{RGB}{0,120,40}
\definecolor{trapcolor}{RGB}{180,80,0}
\definecolor{purpleconcept}{RGB}{120,50,180}
\definecolor{tealhard}{RGB}{0,120,120}

\setbeamercolor{title}{fg=white,bg=redcorrect}
\setbeamercolor{frametitle}{fg=white,bg=redcorrect}
\setbeamercolor{block title}{fg=white,bg=bluecorrect}
\setbeamercolor{block body}{fg=darkgray,bg=lightgray}
\setbeamercolor{normal text}{fg=darkgray}
\usecolortheme{default}

% ====== CUSTOM COMMANDS ======
\newcommand{\correct}[1]{\textcolor{greencheck}{\textbf{\checkmark}} #1}
\newcommand{\keypoint}[1]{\textcolor{deepblue}{\textbf{#1}}}
\newcommand{\hardconcept}[1]{\textcolor{darkred}{\textbf{#1}}}
\newcommand{\examtraps}[1]{\textcolor{trapcolor}{\textbf{⚠ TRAP:}} #1}
\newcommand{\recallbox}[1]{\fcolorbox{deepblue}{blue!5}{\parbox{0.88\textwidth}{\centering \textbf{REMEMBER:} #1}}}
\newcommand{\stepbox}[1]{%
	\begin{center}
		\fcolorbox{darkgray}{white}{%
			\parbox{0.92\textwidth}{\vspace{0.03cm} #1 \vspace{0.03cm}}%
		}
	\end{center}
}
\newcommand{\answerbox}[1]{\fcolorbox{correctgreen}{green!8}{\parbox{0.88\textwidth}{\centering \textcolor{correctgreen}{\textbf{\checkmark\ CORRECT:}} #1}}}
\newcommand{\wronganswer}[1]{\textcolor{redcorrect}{\textbf{✘}} #1}
\newcommand{\trapbox}[1]{\fcolorbox{trapcolor}{orange!8}{\parbox{0.88\textwidth}{\centering \textcolor{trapcolor}{\textbf{⚠ EXAM TRAP:}} #1}}}
\newcommand{\keyrule}[1]{\textcolor{tealhard}{\textbf{🔑}} #1}

\begin{document}
	
	% ================================================================
	% SLIDE 1: TITLE
	% ================================================================
	\begin{frame}
		\centering
		\vspace{0.5cm}
		{\fontsize{18}{22}\selectfont \textbf{CAMS Domain A}}\\[0.15cm]
		{\fontsize{14}{17}\selectfont \textbf{Hard Concepts \& Exam Traps}}\\[0.2cm]
		{\fontsize{9}{11}\selectfont \textit{Advanced Explanations for Difficult Topics}}\\[0.1cm]
		\rule{4cm}{0.5pt}\\[0.1cm]
		{\fontsize{7}{9}\selectfont VASP · Crypto · FATF Recommendations · Sanctions · DORA}\\[0.05cm]
		{\fontsize{7}{9}\selectfont Investment Banking · Wealth Management · PSP · OSINT · KYC}\\[0.05cm]
		{\fontsize{7}{9}\selectfont Trust Companies · Data Governance · Post-Acquisition Integration}
		\vspace{0.3cm}
	\end{frame}
	
	% ================================================================
	% SLIDE 2: VASP - FATF Recommendation 15 DETAILED
	% ================================================================
	\begin{frame}
		\frametitle{\fontsize{11}{13}\selectfont \textbf{VASP \& FATF Recommendation 15 — Complete Breakdown}}
		\begin{block}{\fontsize{7}{9}\selectfont \textbf{What is a VASP?}}
			\scriptsize
			\textbf{Virtual Asset Service Provider} — any natural/legal person conducting \textbf{one or more} of:
			\begin{itemize}
				\item Exchange between virtual assets and fiat currencies
				\item Exchange between one or more forms of virtual assets
				\item Transfer of virtual assets (on behalf of another)
				\item Safekeeping/administration of virtual assets
				\item Participation in financial services related to issuer offers/sale of VAs
			\end{itemize}
		\end{block}
		\vspace{0.02cm}
		\begin{block}{\fontsize{7}{9}\selectfont \textbf{FATF Rec 15 Requirements}}
			\scriptsize
			\begin{itemize}
				\item \textbf{Licensing/Registration:} VASPs must be licensed or registered in all jurisdictions
				\item \textbf{AML/CFT Obligations:} Same as traditional FIs (CDD, EDD, record-keeping, SARs)
				\item \textbf{Travel Rule (Rec 16):} Originator/beneficiary info for all VA transfers
				\item \textbf{Sanctions Screening:} Must screen all wallet addresses against sanctions lists
				\item \textbf{Risk Assessment:} Must assess ML/TF risks of VA products/services
			\end{itemize}
		\end{block}
		\vspace{0.02cm}
		\trapbox{\scriptsize \textbf{Common Exam Trap:} Rec 15 = \textbf{VASP framework}. Rec 16 = \textbf{Travel Rule}. Many candidates confuse these!}
		\vspace{0.02cm}
		\recallbox{\scriptsize Rec 15 is the \textbf{foundational} requirement for VASPs — without it, no VASP can legally operate.}
	\end{frame}
	
	% ================================================================
	% SLIDE 3: CRYPTO AML - HARD CONCEPTS
	% ================================================================
	\begin{frame}
		\frametitle{\fontsize{11}{13}\selectfont \textbf{Cryptocurrency AML — Advanced Concepts}}
		\begin{block}{\fontsize{7}{9}\selectfont \textbf{Key Crypto ML/TF Risks}}
			\scriptsize
			\begin{tabular}{|p{0.28\textwidth}|p{0.68\textwidth}|}
				\hline
				\keypoint{Risk Factor} & \keypoint{Explanation} \\
				\hline
				Pseudonymity & Wallet addresses are not directly linked to real identities \\
				\hline
				Decentralized Exchanges & No KYC required; rapid token swaps obscure source/destination \\
				\hline
				Privacy Coins & Monero, Zcash — transactions are fully anonymous \\
				\hline
				Mixers/Tumblers & Combine funds from multiple sources to obscure origin \\
				\hline
				NFTs & Art can be used for layering via wash trading \\
				\hline
				DeFi Protocols & No central authority; lending/borrowing without KYC \\
				\hline
			\end{tabular}
		\end{block}
		\vspace{0.02cm}
		\begin{block}{\fontsize{7}{9}\selectfont \textbf{Crypto Compliance Toolkit}}
			\scriptsize
			\begin{itemize}
				\item \textbf{Blockchain Analytics:} Chainalysis, Elliptic, CipherTrace — trace on-chain activity
				\item \textbf{Entity Resolution:} Link wallet addresses to real-world entities
				\item \textbf{Wallet Clustering:} Group related wallets to identify entities
				\item \textbf{Real-time Screening:} Screen transactions against sanctions lists
				\item \textbf{Travel Rule Compliance:} Share originator/beneficiary info for transfers > \$1,000
			\end{itemize}
		\end{block}
		\vspace{0.02cm}
		\trapbox{\scriptsize \textbf{Exam Trap:} \wronganswer{Centralized exchanges obscure source/destination} — actually, \textbf{decentralized} exchanges do this. Centralized exchanges require KYC.}
	\end{frame}
	
	% ================================================================
	% SLIDE 4: CRYPTO SANCTIONS SCREENING
	% ================================================================
	\begin{frame}
		\frametitle{\fontsize{11}{13}\selectfont \textbf{Crypto Sanctions Screening — Complete Guide}}
		\begin{block}{\fontsize{7}{9}\selectfont \textbf{Core Requirements}}
			\scriptsize
			\begin{itemize}
				\item \textbf{OFAC Compliance:} Must screen all crypto transactions against SDN List
				\item \textbf{UN Sanctions:} Must screen against UN Security Council sanctions
				\item \textbf{EU Sanctions:} Must comply with EU restrictive measures
				\item \textbf{FATF Rec 7:} Targeted financial sanctions for TF and proliferation financing
				\item \textbf{50\% Rule:} Applies to crypto entities — if 50\%+ owned by sanctioned entity, it's blocked
			\end{itemize}
		\end{block}
		\vspace{0.02cm}
		\begin{block}{\fontsize{7}{9}\selectfont \textbf{Technical Implementation}}
			\scriptsize
			\begin{itemize}
				\item \textbf{On-Chain Screening:} Screen wallet addresses in real-time using blockchain analytics
				\item \textbf{Off-Chain Screening:} Screen counterparty identity info (name, DOB, address)
				\item \textbf{Entity Resolution:} Link on-chain addresses to off-chain risk data
				\item \textbf{Travel Rule:} FATF Rec 16 requires originator/beneficiary info transfer
				\item \textbf{Sanctions Evasion Detection:} Monitor for use of mixers, privacy coins, and cross-chain bridges
			\end{itemize}
		\end{block}
		\vspace{0.02cm}
		\recallbox{\scriptsize \textbf{Key Rule:} Sanctions are \textbf{strict liability} — no intent required. If you process a transaction with a sanctioned entity, you're liable.}
		\vspace{0.02cm}
		\trapbox{\scriptsize \textbf{Exam Trap:} \wronganswer{Only screen fiat transactions} — crypto transactions \textbf{must} be screened too.}
	\end{frame}
	
	% ================================================================
	% SLIDE 5: FATF RECOMMENDATIONS QUICK REFERENCE
	% ================================================================
	\begin{frame}
		\frametitle{\fontsize{11}{13}\selectfont \textbf{FATF Recommendations — Hard-to-Remember Ones}}
		\begin{block}{\fontsize{7}{9}\selectfont \textbf{Critical Recs for CAMS Exam}}
			\scriptsize
			\begin{longtable}{|p{0.12\textwidth}|p{0.40\textwidth}|p{0.44\textwidth}|}
				\hline
				\keypoint{Rec} & \keypoint{Title} & \keypoint{Key Requirement} \\
				\hline
				Rec 1 & Risk-Based Approach & Assess ML/TF risks; apply proportionate measures \\
				\hline
				Rec 6 & Targeted Financial Sanctions (TF) & Freeze assets of designated terrorists \\
				\hline
				Rec 7 & Targeted Financial Sanctions (PF) & Freeze assets of proliferation financiers \\
				\hline
				Rec 10 & Customer Due Diligence & Identify and verify customer identity \\
				\hline
				Rec 12 & PEPs & Apply EDD to foreign PEPs; domestic PEPs risk-based \\
				\hline
				Rec 15 & New Technologies & VASP licensing and AML obligations \\
				\hline
				Rec 16 & Travel Rule & Originator/beneficiary info for wire transfers \\
				\hline
				Rec 17 & Third-Party Reliance & May rely on third parties for CDD (with conditions) \\
				\hline
				Rec 18 & Internal Controls & Group-wide AML programs for all branches/subsidiaries \\
				\hline
				Rec 19 & High-Risk Countries & Apply EDD for FATF black/grey-listed countries \\
				\hline
				Rec 20 & Suspicious Transaction Reporting & File STRs for suspicious activity \\
				\hline
				Rec 22 & DNFBPs & Casinos, real estate, lawyers, accountants, TCSPs \\
				\hline
				Rec 24 & Beneficial Ownership (Legal Persons) & Identify and verify UBOs \\
				\hline
				Rec 25 & Beneficial Ownership (Legal Arrangements) & Trusts: identify settlor, trustee, protector, beneficiaries \\
				\hline
				Rec 35 & Sanctions & Administrative sanctions for non-compliance \\
				\hline
				Rec 36 & Mutual Legal Assistance & International cooperation for investigations \\
				\hline
				Rec 40 & International Cooperation & FIU-to-FIU information sharing \\
				\hline
			\end{longtable}
		\end{block}
		\vspace{0.02cm}
		\recallbox{\scriptsize \textbf{Rec 15 = VASP. Rec 16 = Travel Rule. Rec 17 = Third-party CDD. Rec 24/25 = Beneficial Ownership.}}
	\end{frame}
	
	% ================================================================
	% SLIDE 6: FATF GREY LIST REMEDIATION
	% ================================================================
	\begin{frame}
		\frametitle{\fontsize{11}{13}\selectfont \textbf{FATF Grey List Remediation — Complete Process}}
		\begin{block}{\fontsize{7}{9}\selectfont \textbf{What is the Grey List?}}
			\scriptsize
			Jurisdictions under \textbf{increased monitoring} — committed to resolving strategic AML/CFT deficiencies.
		\end{block}
		\vspace{0.02cm}
		\begin{block}{\fontsize{7}{9}\selectfont \textbf{Remediation Steps (Must-Know Order)}}
			\scriptsize
			\begin{enumerate}
				\item \textbf{Identify Strategic Deficiencies:} FATF publishes specific deficiencies (e.g., weak supervision, BO transparency)
				\item \textbf{Develop Action Plan:} Country submits high-level political commitment to FATF
				\item \textbf{Enact Legislation:} Pass laws to address deficiencies (e.g., AML Act amendments)
				\item \textbf{Strengthen Supervision:} Improve oversight of FIs and DNFBPs
				\item \textbf{Enhance BO Transparency:} Implement BO registers and verification systems
				\item \textbf{Increase Enforcement:} Apply sanctions for non-compliance; prosecute ML cases
				\item \textbf{Show Effectiveness:} Demonstrate that measures are producing results
				\item \textbf{On-Site Visit:} FATF team visits to verify implementation
				\item \textbf{Removal:} FATF votes to remove from grey list
			\end{enumerate}
		\end{block}
		\vspace{0.02cm}
		\begin{block}{\fontsize{7}{9}\selectfont \textbf{Consequences of Grey List Status}}
			\scriptsize
			\begin{itemize}
				\item \textbf{EDD Required:} Rec 19 — enhanced due diligence for transactions involving grey-listed countries
				\item \textbf{Reputational Risk:} Negative perception in financial markets
				\item \textbf{Correspondent Banking:} Banks may terminate relationships
				\item \textbf{Investment:} Reduced foreign direct investment
			\end{itemize}
		\end{block}
		\vspace{0.02cm}
		\trapbox{\scriptsize \textbf{Exam Trap:} \wronganswer{Grey list = black list} — \textbf{NO.} Grey = increased monitoring; Black = call for action. Different consequences!}
	\end{frame}
	
	% ================================================================
	% SLIDE 7: FATF BLACK LIST VS GREY LIST
	% ================================================================
	\begin{frame}
		\frametitle{\fontsize{11}{13}\selectfont \textbf{Black List vs Grey List — Critical Distinction}}
		\begin{block}{\fontsize{7}{9}\selectfont \textbf{Comparison Table}}
			\scriptsize
			\begin{tabular}{|p{0.18\textwidth}|p{0.38\textwidth}|p{0.40\textwidth}|}
				\hline
				\keypoint{Factor} & \keypoint{Black List (Call for Action)} & \keypoint{Grey List (Increased Monitoring)} \\
				\hline
				Definition & Jurisdictions with \textbf{serious} strategic AML/CFT deficiencies & Jurisdictions \textbf{committed} to resolving deficiencies \\
				\hline
				Consequence & \textbf{Strong} countermeasures recommended & \textbf{Enhanced} due diligence (EDD) required \\
				\hline
				FATF Action & Call for action — all members apply countermeasures & Increased monitoring — progress reports required \\
				\hline
				Removal Process & Must address ALL deficiencies; on-site visit & Progress report; may be removed when deficiencies addressed \\
				\hline
				Examples & Iran, North Korea & Nigeria, South Africa, UAE, Turkey \\
				\hline
				Rec Reference & Rec 19 (High-Risk Countries) & Rec 19 (High-Risk Countries) \\
				\hline
			\end{tabular}
		\end{block}
		\vspace{0.02cm}
		\recallbox{\scriptsize \textbf{Black List = Countermeasures (strong sanctions). Grey List = EDD (enhanced due diligence).}}
		\vspace{0.02cm}
		\trapbox{\scriptsize \textbf{Common Error:} \wronganswer{Both lists require the same level of due diligence} — \textbf{NO.} Black list requires \textbf{countermeasures}, not just EDD.}
	\end{frame}
	
	% ================================================================
	% SLIDE 8: INVESTMENT BANKING AML
	% ================================================================
	\begin{frame}
		\frametitle{\fontsize{11}{13}\selectfont \textbf{Investment Banking AML — Hard Concepts}}
		\begin{block}{\fontsize{7}{9}\selectfont \textbf{Key ML Risks in Investment Banking}}
			\scriptsize
			\begin{itemize}
				\item \textbf{Complex Products:} Derivatives, structured products, OTC instruments
				\item \textbf{High-Value Transactions:} Large trades can move significant funds quickly
				\item \textbf{Cross-Border:} Multiple jurisdictions, counterparties, and clearing systems
				\item \textbf{Speed:} Trades execute in milliseconds; monitoring must keep up
				\item \textbf{Anonymity:} Some products can obscure beneficial ownership
			\end{itemize}
		\end{block}
		\vspace{0.02cm}
		\begin{block}{\fontsize{7}{9}\selectfont \textbf{Red Flags Specific to Investment Banking}}
			\scriptsize
			\begin{itemize}
				\item \textbf{Sudden Business Changes:} Client shifts from low-margin to high-return instruments via unfamiliar subsidiary
				\item \textbf{Unusual Trade Patterns:} Wash trading, round-trip transactions, cross-trading
				\item \textbf{Complex Ownership:} Multiple layers, offshore jurisdictions, nominee companies
				\item \textbf{Unusual Counterparties:} Transactions with entities in high-risk jurisdictions
				\item \textbf{Market Manipulation:} Pump-and-dump, spoofing, layering of orders
			\end{itemize}
		\end{block}
		\vspace{0.02cm}
		\begin{block}{\fontsize{7}{9}\selectfont \textbf{Correct Response to Red Flags}}
			\scriptsize
			\begin{itemize}
				\item \textbf{Escalate Internally:} Report to compliance team for investigation
				\item \textbf{Request Documentation:} Obtain beneficial ownership, financial statements, business rationale
				\item \textbf{Conduct EDD:} Enhanced due diligence on counterparties and subsidiaries
				\item \textbf{File SAR/STR:} If suspicious activity is confirmed
			\end{itemize}
		\end{block}
		\vspace{0.02cm}
		\trapbox{\scriptsize \textbf{Exam Trap:} \wronganswer{Immediately decline the client} — first \textbf{escalate and investigate}. De-risking without investigation is not the risk-based approach.}
	\end{frame}
	
	% ================================================================
	% SLIDE 9: WEALTH MANAGEMENT AML
	% ================================================================
	\begin{frame}
		\frametitle{\fontsize{11}{13}\selectfont \textbf{Wealth Management AML — Advanced Concepts}}
		\begin{block}{\fontsize{7}{9}\selectfont \textbf{Core ML Vulnerabilities in Wealth Management}}
			\scriptsize
			\begin{itemize}
				\item \textbf{Discretion:} Relationship managers may have high autonomy, bypassing controls
				\item \textbf{Complex Structures:} Trusts, foundations, offshore companies, nominee arrangements
				\item \textbf{High-Value Clients:} HNWIs with complex financial structures
				\item \textbf{Cross-Border:} Multiple jurisdictions, asset classes, and currencies
				\item \textbf{Confidentiality:} Client expectations of privacy can be abused
			\end{itemize}
		\end{block}
		\vspace{0.02cm}
		\begin{block}{\fontsize{7}{9}\selectfont \textbf{High-Risk Features (Select Three) — CAMS Exam Favorite}}
			\scriptsize
			\begin{enumerate}
				\item \textbf{Minimum Assets Requirement:} Attracts HNWIs with complex structures
				\item \textbf{Offering Numbered Accounts:} Reduces transparency; conceals identity
				\item \textbf{Dedicated Relationship Service with Transaction Discretion:} Enables circumvention of controls
			\end{enumerate}
		\end{block}
		\vspace{0.02cm}
		\begin{block}{\fontsize{7}{9}\selectfont \textbf{Mitigations}}
			\scriptsize
			\begin{itemize}
				\item \textbf{EDD for HNWIs and PEPs:} Source of funds, source of wealth verification
				\item \textbf{Independent Review:} Separate team reviews client transactions
				\item \textbf{Regular Training:} Relationship managers trained on red flags
				\item \textbf{Committee Approval:} High-risk clients require committee-level approval
			\end{itemize}
		\end{block}
		\vspace{0.02cm}
		\trapbox{\scriptsize \textbf{Exam Trap:} \wronganswer{Standardized online account opening} is \textbf{low-risk}, not high-risk. Wealth management risk comes from \textbf{discretion and complexity}.}
	\end{frame}
	
	% ================================================================
	% SLIDE 10: PORTFOLIO MANAGEMENT AML
	% ================================================================
	\begin{frame}
		\frametitle{\fontsize{11}{13}\selectfont \textbf{Portfolio Management AML — Key Concepts}}
		\begin{block}{\fontsize{7}{9}\selectfont \textbf{AML Risks in Portfolio Management}}
			\scriptsize
			\begin{itemize}
				\item \textbf{Discretionary Accounts:} Managers can trade without client approval — risk of ML
				\item \textbf{Complex Products:} Derivatives, structured products, alternative investments
				\item \textbf{Cross-Border:} Investments in multiple jurisdictions
				\item \textbf{High Turnover:} Frequent trading can generate artificial profits/losses
				\item \textbf{Concentration Risk:} Over-concentration in high-risk assets
			\end{itemize}
		\end{block}
		\vspace{0.02cm}
		\begin{block}{\fontsize{7}{9}\selectfont \textbf{Red Flags for Portfolio Managers}}
			\scriptsize
			\begin{itemize}
				\item \textbf{Unusual Investment Strategy:} Client typically low-risk but suddenly invests in high-risk assets
				\item \textbf{Circular Trading:} Client buys and sells same asset through different accounts
				\item \textbf{Wash Trading:} Buying and selling same asset to create artificial volume
				\item \textbf{Unusual Redemptions:} Large, unexpected redemption requests
				\item \textbf{Lack of Investment Objective:} Client has no clear investment strategy
			\end{itemize}
		\end{block}
		\vspace{0.02cm}
		\begin{block}{\fontsize{7}{9}\selectfont \textbf{Compliance Controls}}
			\scriptsize
			\begin{itemize}
				\item \textbf{Pre-Trade Compliance:} Check trades against investment mandates
				\item \textbf{Post-Trade Monitoring:} Review all trades for suspicious patterns
				\item \textbf{Client Risk Assessment:} Regular review of client risk profiles
				\item \textbf{EDD for High-Risk Clients:} Enhanced due diligence for PEPs, HNWIs
			\end{itemize}
		\end{block}
		\vspace{0.02cm}
		\recallbox{\scriptsize \textbf{Key Rule:} Discretionary accounts are \textbf{higher risk} because the manager can trade without client approval.}
	\end{frame}
	
	% ================================================================
	% SLIDE 11: PSP - PAYMENT SERVICE PROVIDERS
	% ================================================================
	\begin{frame}
		\frametitle{\fontsize{11}{13}\selectfont \textbf{PSP AML — Complete Overview}}
		\begin{block}{\fontsize{7}{9}\selectfont \textbf{What is a PSP?}}
			\scriptsize
			\textbf{Payment Service Provider} — entity that enables electronic payments:
			\begin{itemize}
				\item Payment processing (credit/debit cards, direct debits)
				\item Digital wallets (Apple Pay, Google Pay, PayPal)
				\item Money transfer services (Western Union, MoneyGram)
				\item Acquiring services (merchant payment processing)
				\item Payment initiation services (PSD2)
			\end{itemize}
		\end{block}
		\vspace{0.02cm}
		\begin{block}{\fontsize{7}{9}\selectfont \textbf{Key ML Risks for PSPs}}
			\scriptsize
			\begin{itemize}
				\item \textbf{High Volume:} Millions of transactions daily — harder to detect patterns
				\item \textbf{Fast Processing:} Instant settlement — less time for screening
				\item \textbf{Cross-Border:} Multiple jurisdictions with varying AML regimes
				\item \textbf{Agent Networks:} Third-party agents may have weak AML controls
				\item \textbf{Anonymity:} Prepaid cards, digital wallets can be anonymous
			\end{itemize}
		\end{block}
		\vspace{0.02cm}
		\begin{block}{\fontsize{7}{9}\selectfont \textbf{Best Practices for PSPs}}
			\scriptsize
			\begin{itemize}
				\item \textbf{Customize AML Policy:} Per jurisdiction and product
				\item \textbf{EDD for High-Risk Jurisdictions:} Apply enhanced due diligence
				\item \textbf{Real-Time Screening:} Screen all transactions against sanctions lists
				\item \textbf{Agent Oversight:} Monitor agent compliance with AML obligations
				\item \textbf{Transaction Monitoring:} AI/ML to detect patterns in high-volume data
			\end{itemize}
		\end{block}
		\vspace{0.02cm}
		\trapbox{\scriptsize \textbf{Exam Trap:} \wronganswer{Apply home-country standards everywhere} — \textbf{NO.} PSPs must \textbf{customize} per jurisdiction. One-size-fits-all fails.}
	\end{frame}
	
	% ================================================================
	% SLIDE 12: OSINT - OPEN SOURCE INTELLIGENCE
	% ================================================================
	\begin{frame}
		\frametitle{\fontsize{11}{13}\selectfont \textbf{OSINT in AML Investigations — Complete Guide}}
		\begin{block}{\fontsize{7}{9}\selectfont \textbf{What is OSINT?}}
			\scriptsize
			\textbf{Open Source Intelligence} — information collected from publicly available sources:
			\begin{itemize}
				\item Public records (company registers, court records)
				\item Media (news articles, press releases)
				\item Social media (LinkedIn, Facebook, Twitter)
				\item Websites (company websites, government sites)
				\item Databases (sanctions lists, PEP lists, adverse media)
			\end{itemize}
		\end{block}
		\vspace{0.02cm}
		\begin{block}{\fontsize{7}{9}\selectfont \textbf{OSINT in Investigations — When to Use}}
			\scriptsize
			\begin{itemize}
				\item \textbf{Unfamiliar Entity:} Research an unknown counterparty in a sanctioned jurisdiction
				\item \textbf{Beneficial Ownership:} Identify UBOs when documentation is incomplete
				\item \textbf{Adverse Media:} Identify negative news about clients
				\item \textbf{PEP Status:} Confirm PEP status when not flagged by screening
				\item \textbf{Shell Companies:} Identify shell companies and nominee arrangements
			\end{itemize}
		\end{block}
		\vspace{0.02cm}
		\begin{block}{\fontsize{7}{9}\selectfont \textbf{OSINT Best Practices}}
			\scriptsize
			\begin{itemize}
				\item \textbf{Source Verification:} Verify the reliability of sources
				\item \textbf{Documentation:} Record findings and sources for audit trail
				\item \textbf{Privacy Compliance:} Respect data privacy laws (GDPR, CCPA)
				\item \textbf{Tipping Off:} Never disclose investigation to the subject
			\end{itemize}
		\end{block}
		\vspace{0.02cm}
		\recallbox{\scriptsize \textbf{Key Rule:} OSINT is \textbf{critical} for investigating unfamiliar entities. Regulatory advisories are helpful, but \textbf{entity-specific research} is the most useful next step.}
	\end{frame}
	
	% ================================================================
	% SLIDE 13: KYC/CDD - RECOMMENDATION 10
	% ================================================================
	\begin{frame}
		\frametitle{\fontsize{11}{13}\selectfont \textbf{KYC/CDD — FATF Recommendation 10 Detailed}}
		\begin{block}{\fontsize{7}{9}\selectfont \textbf{Core CDD Requirements (Rec 10)}}
			\scriptsize
			\begin{enumerate}
				\item \textbf{Identify Customer:} Name, DOB, address, identification number
				\item \textbf{Verify Identity:} Using reliable, independent source documents/data
				\item \textbf{Identify Beneficial Owner:} Natural person who ultimately owns/controls
				\item \textbf{Verify BO Identity:} Using reliable sources
				\item \textbf{Understand Business Relationship:} Purpose and intended nature
				\item \textbf{Ongoing Due Diligence:} Monitor transactions and update information
			\end{enumerate}
		\end{block}
		\vspace{0.02cm}
		\begin{block}{\fontsize{7}{9}\selectfont \textbf{CDD Timing — Critical Exam Point}}
			\scriptsize
			\begin{itemize}
				\item \textbf{Before} establishing business relationship
				\item \textbf{Before} carrying out occasional transaction (threshold applies)
				\item \textbf{Before} carrying out suspicious transaction
				\item \textbf{When} doubt about previously obtained customer identification data
			\end{itemize}
		\end{block}
		\vspace{0.02cm}
		\begin{block}{\fontsize{7}{9}\selectfont \textbf{Simplified Due Diligence (SDD)}}
			\scriptsize
			\begin{itemize}
				\item \textbf{When Allowed:} Low-risk customers/products/jurisdictions
				\item \textbf{Not Allowed:} If there is suspicion of ML/TF
				\item \textbf{Not Allowed:} If customer is from high-risk jurisdiction
				\item \textbf{Remember:} SDD is \textbf{not zero} due diligence — still must identify and verify
			\end{itemize}
		\end{block}
		\vspace{0.02cm}
		\recallbox{\scriptsize \textbf{Key Rule:} CDD must be done \textbf{before} establishing relationship. Suspicion overrides SDD.}
	\end{frame}
	
	% ================================================================
	% SLIDE 14: ENHANCED DUE DILIGENCE (EDD)
	% ================================================================
	\begin{frame}
		\frametitle{\fontsize{11}{13}\selectfont \textbf{EDD — Enhanced Due Diligence Complete Guide}}
		\begin{block}{\fontsize{7}{9}\selectfont \textbf{When is EDD Required?}}
			\scriptsize
			\begin{itemize}
				\item \textbf{PEPs:} Foreign PEPs \textbf{mandatory}; domestic PEPs \textbf{risk-based}
				\item \textbf{High-Risk Customers:} From FATF black/grey-listed countries
				\item \textbf{High-Risk Products:} Private banking, correspondent banking, crypto
				\item \textbf{High-Risk Jurisdictions:} Countries with weak AML enforcement
				\item \textbf{Complex Ownership:} Trusts, offshore companies, nominee arrangements
				\item \textbf{Suspicion:} If there is suspicion of ML/TF
			\end{itemize}
		\end{block}
		\vspace{0.02cm}
		\begin{block}{\fontsize{7}{9}\selectfont \textbf{EDD Components — MUST KNOW}}
			\scriptsize
			\begin{enumerate}
				\item \textbf{Source of Wealth:} How did the customer acquire their wealth?
				\item \textbf{Source of Funds:} Where did the specific funds come from?
				\item \textbf{Beneficial Ownership:} Identify and verify ALL UBOs (25\%+)
				\item \textbf{Business Purpose:} Understand the reason for the relationship
				\item \textbf{Ongoing Monitoring:} Enhanced transaction monitoring
				\item \textbf{Senior Approval:} Senior management approval required
				\item \textbf{Regular Reviews:} More frequent reviews (e.g., annually vs. 3 years)
			\end{enumerate}
		\end{block}
		\vspace{0.02cm}
		\begin{block}{\fontsize{7}{9}\selectfont \textbf{SOW vs SOF — Critical Distinction}}
			\scriptsize
			\begin{itemize}
				\item \textbf{Source of Wealth:} \textit{How} did you get your money? (e.g., inheritance, business profits)
				\item \textbf{Source of Funds:} \textit{Where} did this specific money come from? (e.g., sale of property, business sale)
			\end{itemize}
		\end{block}
		\vspace{0.02cm}
		\trapbox{\scriptsize \textbf{Exam Trap:} \wronganswer{SOW and SOF are the same} — \textbf{NO.} SOW = overall wealth; SOF = specific funds for the transaction.}
	\end{frame}
	
	% ================================================================
	% SLIDE 15: TRUST AND COMPANY SERVICE PROVIDERS (TCSP)
	% ================================================================
	\begin{frame}
		\frametitle{\fontsize{11}{13}\selectfont \textbf{TCSP — Trust \& Company Service Providers (Rec 22)}}
		\begin{block}{\fontsize{7}{9}\selectfont \textbf{What is a TCSP? (DNFBP)}}
			\scriptsize
			\textbf{Trust and Company Service Providers} are \textbf{DNFBPs} under FATF Rec 22:
			\begin{itemize}
				\item Incorporation of companies
				\item Administration of trusts and companies
				\item Acting as company secretary
				\item Providing registered office address
				\item Acting as director/nominee shareholder
				\item Acting as trustee or protector of a trust
			\end{itemize}
		\end{block}
		\vspace{0.02cm}
		\begin{block}{\fontsize{7}{9}\selectfont \textbf{TCSP AML Obligations}}
			\scriptsize
			\begin{itemize}
				\item \textbf{CDD:} Identify and verify customers and beneficial owners (Rec 10)
				\item \textbf{BO Identification:} Identify natural persons who own/control (Rec 24/25)
				\item \textbf{Record-Keeping:} Maintain records for 5+ years
				\item \textbf{STR Filing:} File suspicious transaction reports
				\item \textbf{Internal Controls:} Implement AML program
				\item \textbf{Sanctions Screening:} Screen against sanctions lists
			\end{itemize}
		\end{block}
		\vspace{0.02cm}
		\begin{block}{\fontsize{7}{9}\selectfont \textbf{High-Risk Client Handling — CAMS Exam Favorite}}
			\scriptsize
			\textbf{Scenario:} Client requests confidentiality for all shareholders. Connection to high-risk industries in NRA.
			\par \textbf{Correct Response:} \textbf{Decline immediately} unless full BO information and risk justifications are disclosed and verified.
		\end{block}
		\vspace{0.02cm}
		\trapbox{\scriptsize \textbf{Exam Trap:} \wronganswer{Refer to local counsel and detach from compliance responsibility} — \textbf{NO.} TCSP retains compliance responsibility. Cannot outsource AML obligations.}
	\end{frame}
	
	% ================================================================
	% SLIDE 16: DATA GOVERNANCE - BCBS 239
	% ================================================================
	\begin{frame}
		\frametitle{\fontsize{11}{13}\selectfont \textbf{Data Governance — BCBS 239 \& AML}}
		\begin{block}{\fontsize{7}{9}\selectfont \textbf{BCBS 239 — Principles for Effective Risk Data Aggregation}}
			\scriptsize
			\textbf{14 Principles} — banks must have:
			\begin{itemize}
				\item \textbf{Accuracy \& Integrity:} Data must be correct and reliable
				\item \textbf{Completeness:} All data must be captured
				\item \textbf{Timeliness:} Data must be available when needed
				\item \textbf{Adaptability:} Systems must adapt to changing needs
			\end{itemize}
		\end{block}
		\vspace{0.02cm}
		\begin{block}{\fontsize{7}{9}\selectfont \textbf{AML Data Governance Challenges}}
			\scriptsize
			\begin{itemize}
				\item \textbf{Missing Identifiers:} No DOB, no tax ID, no passport number
				\item \textbf{Inconsistent Formats:} "US" vs "USA" vs "840" vs "United States"
				\item \textbf{Duplicate Profiles:} Same customer appears multiple times
				\item \textbf{Incomplete Data:} Missing fields in transaction monitoring
				\item \textbf{Data Silos:} Different systems don't share data
			\end{itemize}
		\end{block}
		\vspace{0.02cm}
		\begin{block}{\fontsize{7}{9}\selectfont \textbf{Data Cleansing Process — CAMS Exam Favorite}}
			\scriptsize
			\begin{enumerate}
				\item \textbf{Removal of Duplicate Entries:} Merge duplicate customer profiles
				\item \textbf{Validation of Data Accuracy:} Check formats and values
				\item \textbf{Standardization:} Apply consistent formats (e.g., country codes)
				\item \textbf{Enrichment:} Add missing data from external sources
			\end{enumerate}
		\end{block}
		\vspace{0.02cm}
		\recallbox{\scriptsize \textbf{Key Rule:} Data cleansing is \textbf{pre-requisite} for effective AML controls. Without clean data, monitoring alerts are unreliable.}
	\end{frame}
	
	% ================================================================
	% SLIDE 17: DATA LINEAGE - TRANSPARENCY
	% ================================================================
	\begin{frame}
		\frametitle{\fontsize{11}{13}\selectfont \textbf{Data Lineage — Transparency \& Auditability}}
		\begin{block}{\fontsize{7}{9}\selectfont \textbf{What is Data Lineage?}}
			\scriptsize
			\textbf{Data lineage} is the process of tracking the \textbf{flow of data} from its source to its destination:
			\begin{itemize}
				\item Where did the data come from?
				\item How was it transformed?
				\item Where is it stored?
				\item Who accessed it?
				\item When was it modified?
			\end{itemize}
		\end{block}
		\vspace{0.02cm}
		\begin{block}{\fontsize{7}{9}\selectfont \textbf{Data Lineage for AML — CAMS Exam Favorite}}
			\scriptsize
			What supports transparency and auditability through data lineage? (Select Two.)
			\begin{enumerate}
				\item \textbf{Maintaining history of data transformations}
				\item \textbf{Capturing source of each data field}
			\end{enumerate}
		\end{block}
		\vspace{0.02cm}
		\begin{block}{\fontsize{7}{9}\selectfont \textbf{Why Data Lineage Matters in AML}}
			\scriptsize
			\begin{itemize}
				\item \textbf{Regulatory Audit:} Regulators want to see where data came from
				\item \textbf{Model Validation:} Need to trace data inputs to models
				\item \textbf{Error Correction:} Trace errors back to source
				\item \textbf{SAR Filing:} Need to verify accuracy of reported data
				\item \textbf{Data Quality:} Identify and correct data quality issues
			\end{itemize}
		\end{block}
		\vspace{0.02cm}
		\trapbox{\scriptsize \textbf{Exam Trap:} \wronganswer{Archiving analyst notes} is case documentation, \textbf{not} data lineage. \wronganswer{Deleting outdated data} undermines auditability.}
	\end{frame}
	
	% ================================================================
	% SLIDE 18: POST-ACQUISITION AML INTEGRATION
	% ================================================================
	\begin{frame}
		\frametitle{\fontsize{11}{13}\selectfont \textbf{Post-Acquisition AML Integration — Complete Process}}
		\begin{block}{\fontsize{7}{9}\selectfont \textbf{Scenario — CAMS Exam Favorite}}
			\scriptsize
			During post-acquisition integration, compliance team reviews merged customer databases and identifies:
			\begin{itemize}
				\item Multiple entries with inconsistent address formats
				\item Duplicate customer profiles
				\item Missing identifiers (DOB, tax ID)
				\item Inconsistent country codes
			\end{itemize}
		\end{block}
		\vspace{0.02cm}
		\begin{block}{\fontsize{7}{9}\selectfont \textbf{Correct Priority — Step-by-Step}}
			\scriptsize
			\begin{enumerate}
				\item \textbf{Data Cleansing:} Target missing and duplicated records \textbf{immediately}
				\item \textbf{Data Coverage \& Gap Assessment:} Identify what data is missing
				\item \textbf{Standardization:} Harmonize country codes and data formats
				\item \textbf{System Integration:} Align transaction monitoring scenarios
				\item \textbf{Policy Alignment:} Harmonize AML policies across merged entities
				\item \textbf{Training:} Train staff on new systems and procedures
			\end{enumerate}
		\end{block}
		\vspace{0.02cm}
		\begin{block}{\fontsize{7}{9}\selectfont \textbf{Why Data Cleansing First?}}
			\scriptsize
			\begin{itemize}
				\item Without clean data, you cannot:
				\begin{itemize}
					\item Screen accurately against sanctions
					\item Monitor transactions effectively
					\item File accurate SARs
					\item Assess risk properly
				\end{itemize}
			\end{itemize}
		\end{block}
		\vspace{0.02cm}
		\recallbox{\scriptsize \textbf{Key Rule:} Data cleansing is the \textbf{first} priority in post-acquisition integration. Clean data is the foundation of AML controls.}
	\end{frame}
	
	% ================================================================
	% SLIDE 19: SANCTIONS SCREENING - COMPLETE
	% ================================================================
	\begin{frame}
		\frametitle{\fontsize{11}{13}\selectfont \textbf{Sanctions Screening — Complete Guide}}
		\begin{block}{\fontsize{7}{9}\selectfont \textbf{What is Sanctions Screening?}}
			\scriptsize
			Process of checking customers, transactions, and counterparties against \textbf{sanctions lists}:
			\begin{itemize}
				\item OFAC SDN List (US)
				\item UN Security Council Sanctions (UN)
				\item EU Restrictive Measures (EU)
				\item HM Treasury Sanctions List (UK)
				\item FATF High-Risk Jurisdictions (Black/Grey Lists)
			\end{itemize}
		\end{block}
		\vspace{0.02cm}
		\begin{block}{\fontsize{7}{9}\selectfont \textbf{Sanctions Screening Process}}
			\scriptsize
			\begin{itemize}
				\item \textbf{Real-Time Screening:} At point of onboarding, transactions, payments
				\item \textbf{Batch Screening:} Periodic re-screening of entire customer database
				\item \textbf{Name Screening:} Match against names on sanctions lists
				\item \textbf{Transaction Screening:} Match against sanctions list entries
				\item \textbf{Blocking:} Must block/restrict transactions with sanctioned entities
			\end{itemize}
		\end{block}
		\vspace{0.02cm}
		\begin{block}{\fontsize{7}{9}\selectfont \textbf{OFAC 50\% Rule — MUST KNOW}}
			\scriptsize
			If a sanctioned entity owns \textbf{50\% or more} of another entity, that entity is \textbf{also blocked} — even if not on SDN list.
		\end{block}
		\vspace{0.02cm}
		\begin{block}{\fontsize{7}{9}\selectfont \textbf{External Sanctions Sources — CAMS Exam}}
			\scriptsize
			Which are \textbf{external} sources? (Select Two.)
			\begin{enumerate}
				\item \textbf{FATF blacklists and UN Security Council lists}
				\item \textbf{State Department embargo registries}
			\end{enumerate}
		\end{block}
		\vspace{0.02cm}
		\recallbox{\scriptsize \textbf{Key Rule:} Sanctions are \textbf{strict liability}. No intent required — if you process a transaction with a sanctioned entity, you're liable.}
	\end{frame}
	
	% ================================================================
	% SLIDE 20: TRANSACTION MONITORING - ADVANCED
	% ================================================================
	\begin{frame}
		\frametitle{\fontsize{11}{13}\selectfont \textbf{Transaction Monitoring — Advanced Concepts}}
		\begin{block}{\fontsize{7}{9}\selectfont \textbf{Transaction Monitoring System Components}}
			\scriptsize
			\begin{itemize}
				\item \textbf{Data Ingestion:} Collect transaction data from multiple systems
				\item \textbf{Scenarios/Rules:} Pre-defined rules to detect suspicious patterns
				\item \textbf{Scoring:} Assign risk scores to transactions/customers
				\item \textbf{Alert Generation:} Create alerts for suspicious activity
				\item \textbf{Case Management:} Manage investigation workflow
				\item \textbf{Reporting:} Generate SARs/STRs
			\end{itemize}
		\end{block}
		\vspace{0.02cm}
		\begin{block}{\fontsize{7}{9}\selectfont \textbf{Rules-Based vs AI/ML Monitoring}}
			\scriptsize
			\begin{tabular}{|p{0.30\textwidth}|p{0.32\textwidth}|p{0.34\textwidth}|}
				\hline
				\keypoint{Factor} & \keypoint{Rules-Based} & \keypoint{AI/ML} \\
				\hline
				False Positives & High & Lower (with tuning) \\
				\hline
				Adaptability & Static; requires manual updates & Dynamic; learns from data \\
				\hline
				Complex Patterns & Limited detection & Advanced pattern detection \\
				\hline
				Transparency & High (rules are clear) & Lower (black box concern) \\
				\hline
				Regulatory Acceptance & Well-established & Requires validation \\
				\hline
			\end{tabular}
		\end{block}
		\vspace{0.02cm}
		\begin{block}{\fontsize{7}{9}\selectfont \textbf{Typical Scenarios (Red Flags)}}
			\scriptsize
			\begin{itemize}
				\item Structuring: Multiple small deposits below reporting threshold
				\item Round-Trip Transactions: Funds leave and return to same account
				\item Rapid Movement: Funds move quickly through multiple accounts
				\item Unusual Geographic Patterns: Transactions to high-risk jurisdictions
				\item Mismatched Profile: Transaction pattern doesn't match KYC profile
			\end{itemize}
		\end{block}
		\vspace{0.02cm}
		\recallbox{\scriptsize \textbf{Key Rule:} Transaction monitoring is about \textbf{detecting unusual activity} — not just large transactions.}
	\end{frame}
	
	% ================================================================
	% SLIDE 21: GAMBLING AML
	% ================================================================
	\begin{frame}
		\frametitle{\fontsize{11}{13}\selectfont \textbf{Gambling AML — Complete Overview}}
		\begin{block}{\fontsize{7}{9}\selectfont \textbf{Why Gambling is High-Risk for ML}}
			\scriptsize
			\begin{itemize}
				\item \textbf{Cash-Intensive:} Large volumes of cash transactions
				\item \textbf{Anonymity:} Chips can be purchased anonymously
				\item \textbf{Rapid Movement:} Funds move quickly in/out
				\item \textbf{Cross-Border:} Online gambling operates globally
				\item \textbf{Layering:} Chips can be used to layer illicit funds
				\item \textbf{Integration:} Winnings appear legitimate
			\end{itemize}
		\end{block}
		\vspace{0.02cm}
		\begin{block}{\fontsize{7}{9}\selectfont \textbf{Casino AML Controls — CAMS Exam Favorite}}
			\scriptsize
			Which are effective controls? (Select Two.)
			\begin{enumerate}
				\item \textbf{Enforcing chip redemption policies tied to minimum gaming thresholds}
				\item \textbf{Tracking player activity through registered loyalty programs}
			\end{enumerate}
		\end{block}
		\vspace{0.02cm}
		\begin{block}{\fontsize{7}{9}\selectfont \textbf{Online Gambling Risks}}
			\scriptsize
			\begin{itemize}
				\item \textbf{Cryptocurrency:} Used to fund gambling accounts
				\item \textbf{E-Wallets:} Rapid movement of funds
				\item \textbf{Multiple Accounts:} Players use multiple accounts
				\item \textbf{Bonus Abuse:} Using bonuses to launder funds
				\item \textbf{P2P Transfers:} Players transfer funds between accounts
			\end{itemize}
		\end{block}
		\vspace{0.02cm}
		\begin{block}{\fontsize{7}{9}\selectfont \textbf{Countermeasures}}
			\scriptsize
			\begin{itemize}
				\item \textbf{CDD:} Identify and verify customers
				\item \textbf{Transaction Monitoring:} Track deposits, withdrawals, and gameplay
				\item \textbf{Threshold Reporting:} Report transactions above thresholds
				\item \textbf{EDD:} Enhanced due diligence for high-risk customers
				\item \textbf{Sanctions Screening:} Screen customers against sanctions lists
			\end{itemize}
		\end{block}
		\vspace{0.02cm}
		\recallbox{\scriptsize \textbf{Key Rule:} Casinos are \textbf{DNFBPs} under FATF Rec 22 and must have AML programs.}
	\end{frame}
	
	% ================================================================
	% SLIDE 22: ASSET MANAGER AML
	% ================================================================
	\begin{frame}
		\frametitle{\fontsize{11}{13}\selectfont \textbf{Asset Manager AML — Key Concepts}}
		\begin{block}{\fontsize{7}{9}\selectfont \textbf{AML Risks for Asset Managers}}
			\scriptsize
			\begin{itemize}
				\item \textbf{Complex Products:} Alternative investments, hedge funds, private equity
				\item \textbf{Cross-Border:} Multiple jurisdictions and counterparties
				\item \textbf{High-Value:} Large investments
				\item \textbf{Anonymity:} Some products can obscure beneficial ownership
				\item \textbf{Third-Party Assets:} Managing assets for third parties
			\end{itemize}
		\end{block}
		\vspace{0.02cm}
		\begin{block}{\fontsize{7}{9}\selectfont \textbf{Red Flags for Asset Managers}}
			\scriptsize
			\begin{itemize}
				\item \textbf{Unusual Investment Strategy:} Mismatch with client profile
				\item \textbf{Unusual Redemptions:} Large, unexpected redemption requests
				\item \textbf{Circular Trading:} Client buys/sells same asset
				\item \textbf{Concentration Risk:} Over-concentration in high-risk assets
				\item \textbf{Source of Funds:} Inability to verify source of funds
			\end{itemize}
		\end{block}
		\vspace{0.02cm}
		\begin{block}{\fontsize{7}{9}\selectfont \textbf{Compliance Controls}}
			\scriptsize
			\begin{itemize}
				\item \textbf{CDD/EDD:} Identify and verify customers and BOs
				\item \textbf{Transaction Monitoring:} Monitor all trades and transactions
				\item \textbf{Sanctions Screening:} Screen against sanctions lists
				\item \textbf{PEP Screening:} Identify and monitor PEPs
				\item \textbf{Adverse Media:} Monitor for negative news about clients
				\item \textbf{Independent Review:} Review by compliance team
			\end{itemize}
		\end{block}
		\vspace{0.02cm}
		\recallbox{\scriptsize \textbf{Key Rule:} Asset managers are \textbf{financial institutions} under FATF and must have comprehensive AML programs.}
	\end{frame}
	
	% ================================================================
	% SLIDE 23: FRAUD - AML INTERSECTION
	% ================================================================
	\begin{frame}
		\frametitle{\fontsize{11}{13}\selectfont \textbf{Fraud \& AML — Critical Intersection}}
		\begin{block}{\fontsize{7}{9}\selectfont \textbf{Fraud as a Predicate Offense}}
			\scriptsize
			Fraud is a \textbf{predicate offense} for money laundering:
			\begin{itemize}
				\item Proceeds of fraud must be laundered
				\item Fraudsters use laundering techniques to hide proceeds
				\item AML controls can detect fraud
			\end{itemize}
		\end{block}
		\vspace{0.02cm}
		\begin{block}{\fontsize{7}{9}\selectfont \textbf{Fraud Types Relevant to AML}}
			\scriptsize
			\begin{itemize}
				\item \textbf{Synthetic Identity Fraud:} Creating fake identities using real/partially real info
				\item \textbf{Account Takeover:} Criminal takes over legitimate account
				\item \textbf{Payment Fraud:} Fraudulent payments and transfers
				\item \textbf{Application Fraud:} Fraudulent loan/credit card applications
				\item \textbf{Phishing:} Stealing credentials via deceptive emails
				\item \textbf{Ransomware:} Extorting payments via crypto
			\end{itemize}
		\end{block}
		\vspace{0.02cm}
		\begin{block}{\fontsize{7}{9}\selectfont \textbf{Corrective Action for Fraud — CAMS Exam}}
			\scriptsize
			\textbf{Scenario:} New credit product targeted at low-income customers. Firm identifies trend of synthetic identity fraud.
			\par \textbf{Correct Action:} \textbf{Suspend the product} while updating risk assessment and onboarding controls.
		\end{block}
		\vspace{0.02cm}
		\recallbox{\scriptsize \textbf{Key Rule:} Fraud and AML are \textbf{integrated}. AML controls detect fraud, and fraud controls detect AML.}
	\end{frame}
	
	% ================================================================
	% SLIDE 24: DORA - DIGITAL OPERATIONAL RESILIENCE ACT
	% ================================================================
	\begin{frame}
		\frametitle{\fontsize{11}{13}\selectfont \textbf{DORA — Digital Operational Resilience Act (Complete)}}
		\begin{block}{\fontsize{7}{9}\selectfont \textbf{What is DORA?}}
			\scriptsize
			\textbf{Regulation (EU) 2022/2554} — applies to all financial entities operating in the EU:
			\begin{itemize}
				\item Banks, investment firms, credit institutions
				\item Crypto-asset service providers
				\item Third-party ICT providers
				\item Payment institutions, e-money institutions
				\item Insurance and reinsurance companies
			\end{itemize}
		\end{block}
		\vspace{0.02cm}
		\begin{block}{\fontsize{7}{9}\selectfont \textbf{Key Regulatory Expectation — CAMS Exam Favorite}}
			\scriptsize
			\begin{center}
				\textbf{Centralizing oversight} of ICT risk management across all member-state branches.
			\end{center}
		\end{block}
		\vspace{0.02cm}
		\begin{block}{\fontsize{7}{9}\selectfont \textbf{DORA Requirements}}
			\scriptsize
			\begin{itemize}
				\item \textbf{ICT Risk Governance Framework:} Unified across all EU operations
				\item \textbf{ICT Risk Assessments:} Identify and assess ICT risks
				\item \textbf{Incident Reporting:} Report ICT incidents to regulators
				\item \textbf{Third-Party ICT Provider Risk Management:} Monitor and manage vendor risks
				\item \textbf{Resilience Testing:} Regular testing of ICT systems
				\item \textbf{Information Sharing:} Share threat intelligence with peers
			\end{itemize}
		\end{block}
		\vspace{0.02cm}
		\trapbox{\scriptsize \textbf{Exam Trap:} \wronganswer{Relying solely on national regulators' cybersecurity requirements} is \textbf{insufficient}. DORA mandates centralized, pan-European oversight.}
	\end{frame}
	
	% ================================================================
	% SLIDE 25: RECOMMENDATION 17 - THIRD-PARTY RELIANCE
	% ================================================================
	\begin{frame}
		\frametitle{\fontsize{11}{13}\selectfont \textbf{Rec 17 — Third-Party Reliance for CDD}}
		\begin{block}{\fontsize{7}{9}\selectfont \textbf{What is Third-Party Reliance?}}
			\scriptsize
			Financial institutions may \textbf{rely on third parties} to perform CDD (Rec 17):
			\begin{itemize}
				\item Another financial institution (in same group)
				\item Third-party service provider (if regulated)
				\item Professional advisor (lawyer, accountant)
				\item Trust/company service provider (TCSP)
			\end{itemize}
		\end{block}
		\vspace{0.02cm}
		\begin{block}{\fontsize{7}{9}\selectfont \textbf{Conditions for Reliance — MUST KNOW}}
			\scriptsize
			\begin{enumerate}
				\item Third party must be \textbf{subject to FATF-compliant AML/CFT requirements}
				\item Third party must \textbf{apply CDD requirements} equivalent to Rec 10
				\item Third party must \textbf{obtain and verify} customer identification
				\item Third party must \textbf{provide the information promptly} upon request
				\item The \textbf{ultimate responsibility} remains with the relying institution
				\item The relying institution must be \textbf{satisfied} that the third party will comply
			\end{enumerate}
		\end{block}
		\vspace{0.02cm}
		\begin{block}{\fontsize{7}{9}\selectfont \textbf{Critical Exam Points}}
			\scriptsize
			\begin{itemize}
				\item \textbf{Liability:} The relying institution remains \textbf{fully liable} for CDD
				\item \textbf{Data Protection:} Must comply with data protection laws (GDPR)
				\item \textbf{Record-Keeping:} Must maintain records of reliance
				\item \textbf{Risk-Based:} Can rely for low-risk; may need EDD for high-risk
			\end{itemize}
		\end{block}
		\vspace{0.02cm}
		\recallbox{\scriptsize \textbf{Key Rule:} You can \textbf{rely} on third parties for CDD, but you \textbf{cannot outsource responsibility}. You remain liable.}
	\end{frame}
	
	% ================================================================
	% SLIDE 26: CROSS-BORDER DATA PRIVACY
	% ================================================================
	\begin{frame}
		\frametitle{\fontsize{11}{13}\selectfont \textbf{Cross-Border Data Privacy — AML Compliance}}
		\begin{block}{\fontsize{7}{9}\selectfont \textbf{Key Data Privacy Regulations}}
			\scriptsize
			\begin{itemize}
				\item \textbf{GDPR (EU):} General Data Protection Regulation
				\item \textbf{CCPA (California):} California Consumer Privacy Act
				\item \textbf{PDPA (Singapore):} Personal Data Protection Act
				\item \textbf{LGPD (Brazil):} Lei Geral de Proteção de Dados
			\end{itemize}
		\end{block}
		\vspace{0.02cm}
		\begin{block}{\fontsize{7}{9}\selectfont \textbf{AML vs Data Privacy Conflict}}
			\scriptsize
			\begin{tabular}{|p{0.40\textwidth}|p{0.56\textwidth}|}
				\hline
				\keypoint{AML Requirement} & \keypoint{Data Privacy Concern} \\
				\hline
				Collect customer data & Data minimization principle \\
				\hline
				Share with regulators & Data transfer restrictions \\
				\hline
				Retain records for 5+ years & Storage limitation principle \\
				\hline
				Share within group & Data sharing restrictions \\
				\hline
				Screen against watchlists & Processing sensitive data \\
				\hline
			\end{tabular}
		\end{block}
		\vspace{0.02cm}
		\begin{block}{\fontsize{7}{9}\selectfont \textbf{Correct Approach — CAMS Exam}}
			\scriptsize
			\begin{enumerate}
				\item \textbf{Jurisdiction-Specific Procedures:} Create procedures that comply with local laws
				\item \textbf{Standard Template with Adjustable Sections:} Core sections consistent; adjustable for local laws
				\item \textbf{Legal Basis:} Ensure legal basis for processing (e.g., legal obligation)
				\item \textbf{Data Transfer Mechanisms:} Use SCCs, BCRs, or other transfer mechanisms
				\item \textbf{Data Protection Impact Assessment:} Conduct DPIA for high-risk processing
			\end{enumerate}
		\end{block}
		\vspace{0.02cm}
		\recallbox{\scriptsize \textbf{Key Rule:} AML compliance \textbf{does not override} data privacy. Must comply with both.}
	\end{frame}
	
	% ================================================================
	% SLIDE 27: FREE-TRADE ZONES (FTZ) AML
	% ================================================================
	\begin{frame}
		\frametitle{\fontsize{11}{13}\selectfont \textbf{Free-Trade Zones (FTZ) — AML Vulnerabilities}}
		\begin{block}{\fontsize{7}{9}\selectfont \textbf{What are FTZs?}}
			\scriptsize
			Special economic zones with \textbf{reduced customs and regulatory requirements}:
			\begin{itemize}
				\item Duty-free import/export
				\item Reduced inspections
				\item Limited customs authority
				\item Special tax regimes
			\end{itemize}
		\end{block}
		\vspace{0.02cm}
		\begin{block}{\fontsize{7}{9}\selectfont \textbf{FTZ Vulnerabilities — CAMS Exam Favorite}}
			\scriptsize
			Which are typical FTZ vulnerabilities? (Select Two.)
			\begin{enumerate}
				\item \textbf{Limited customs authority over transitory goods}
				\item \textbf{Evasions of customs payments}
			\end{enumerate}
		\end{block}
		\vspace{0.02cm}
		\begin{block}{\fontsize{7}{9}\selectfont \textbf{FTZ ML Techniques}}
			\scriptsize
			\begin{itemize}
				\item \textbf{Mis-Invoicing:} Over/under-invoice goods to move value
				\item \textbf{Phantom Shipments:} Invoicing for goods that don't exist
				\item \textbf{Customs Evasion:} Declaring goods for re-export, then selling locally
				\item \textbf{Shell Companies:} Using FTZ shell companies to hide ownership
				\item \textbf{Commingling:} Mixing legal and illegal goods
			\end{itemize}
		\end{block}
		\vspace{0.02cm}
		\begin{block}{\fontsize{7}{9}\selectfont \textbf{Mitigations}}
			\scriptsize
			\begin{itemize}
				\item \textbf{EDD:} Enhanced due diligence for FTZ-related clients
				\item \textbf{Transaction Monitoring:} Monitor FTZ-related transactions
				\item \textbf{Document Verification:} Verify shipping and customs documents
				\item \textbf{Sanctions Screening:} Screen FTZ counterparties against sanctions
			\end{itemize}
		\end{block}
		\vspace{0.02cm}
		\trapbox{\scriptsize \textbf{Exam Trap:} \wronganswer{Over-invoicing} is a TBML technique, but FTZ-specific risk in scenario = \textbf{customs payment evasion}.}
	\end{frame}
	
	% ================================================================
	% SLIDE 28: TBML - TRADE-BASED MONEY LAUNDERING
	% ================================================================
	\begin{frame}
		\frametitle{\fontsize{11}{13}\selectfont \textbf{TBML — Trade-Based Money Laundering}}
		\begin{block}{\fontsize{7}{9}\selectfont \textbf{What is TBML?}}
			\scriptsize
			\textbf{Trade-Based Money Laundering} — using international trade to launder money:
			\begin{itemize}
				\item Mis-invoicing (over/under-invoicing)
				\item Phantom shipments
				\item Multiple invoicing
				\item Over/under-shipment
				\item Using shell companies
			\end{itemize}
		\end{block}
		\vspace{0.02cm}
		\begin{block}{\fontsize{7}{9}\selectfont \textbf{Key TBML Techniques — CAMS Exam}}
			\scriptsize
			\begin{tabular}{|p{0.30\textwidth}|p{0.35\textwidth}|p{0.31\textwidth}|}
				\hline
				\keypoint{Technique} & \keypoint{Description} & \keypoint{Beneficiary} \\
				\hline
				Over-Invoicing & Invoice price higher than market & Seller (exporter) \\
				\hline
				Under-Invoicing & Invoice price lower than market & Buyer (importer) \\
				\hline
				Phantom Shipment & Invoice for goods not shipped & Both parties \\
				\hline
				Multiple Invoicing & Multiple invoices for same goods & Seller (exporter) \\
				\hline
				Over-Shipment & Ship more than invoiced & Buyer (importer) \\
				\hline
			\end{tabular}
		\end{block}
		\vspace{0.02cm}
		\begin{block}{\fontsize{7}{9}\selectfont \textbf{TBML Red Flags}}
			\scriptsize
			\begin{itemize}
				\item \textbf{Pricing:} Significantly above/below market rates
				\item \textbf{Documentation:} Inconsistencies in shipping documents
				\item \textbf{Counterparties:} Transactions with high-risk jurisdictions
				\item \textbf{Routes:} Unusual shipping routes
				\item \textbf{Payment:} Third-party payments
			\end{itemize}
		\end{block}
		\vspace{0.02cm}
		\recallbox{\scriptsize \textbf{Key Rule:} Over-invoicing = value to \textbf{seller}. Under-invoicing = value to \textbf{buyer}.}
	\end{frame}
	
	% ================================================================
	% SLIDE 29: STRUCTURING - DETAILED
	% ================================================================
	\begin{frame}
		\frametitle{\fontsize{11}{13}\selectfont \textbf{Structuring (Smurfing) — Complete Guide}}
		\begin{block}{\fontsize{7}{9}\selectfont \textbf{What is Structuring?}}
			\scriptsize
			Breaking up large transactions into \textbf{smaller amounts} to avoid reporting thresholds.
			\par \textbf{Thresholds:}
			\begin{itemize}
				\item US: \$10,000 (CTR filing)
				\item EU: €10,000 (Cash controls)
				\item UK: £10,000 (SAR threshold)
			\end{itemize}
		\end{block}
		\vspace{0.02cm}
		\begin{block}{\fontsize{7}{9}\selectfont \textbf{Key Fact — CAMS Exam Favorite}}
			\scriptsize
			\begin{center}
				\textbf{Structuring is illegal even if the money is legitimate.}
			\end{center}
			The crime is \textbf{evading reporting requirements} — not the source of funds.
		\end{block}
		\vspace{0.02cm}
		\begin{block}{\fontsize{7}{9}\selectfont \textbf{Structuring Red Flags}}
			\scriptsize
			\begin{itemize}
				\item Multiple small cash deposits just below the reporting threshold
				\item Deposits made by different people into the same account
				\item Deposits made at different branches/locations
				\item Deposits made over several days/weeks
				\item Sudden international wire transfer after structuring
			\end{itemize}
		\end{block}
		\vspace{0.02cm}
		\begin{block}{\fontsize{7}{9}\selectfont \textbf{Analyst Action — CAMS Exam}}
			\scriptsize
			\begin{itemize}
				\item Identify structuring pattern
				\item \textbf{File SAR} (suspicious activity report)
				\item \textbf{DO NOT tip off} the customer
				\item Escalate to Level 2 for investigation
			\end{itemize}
		\end{block}
		\vspace{0.02cm}
		\recallbox{\scriptsize \textbf{Key Rule:} Structuring is illegal \textbf{regardless} of whether the money is clean or dirty.}
	\end{frame}
	
	% ================================================================
	% SLIDE 30: LAYERING - DETAILED
	% ================================================================
	\begin{frame}
		\frametitle{\fontsize{11}{13}\selectfont \textbf{Layering — Money Laundering Stage 2}}
		\begin{block}{\fontsize{7}{9}\selectfont \textbf{What is Layering?}}
			\scriptsize
			\textbf{Second stage} of money laundering — \textbf{concealing the origin} of illicit funds:
			\begin{itemize}
				\item Moving funds through multiple accounts
				\item Converting funds to different assets
				\item Using complex transactions
				\item Creating false paper trails
				\item Distance from source
			\end{itemize}
		\end{block}
		\vspace{0.02cm}
		\begin{block}{\fontsize{7}{9}\selectfont \textbf{Layering Techniques — CAMS Exam}}
			\scriptsize
			\begin{itemize}
				\item \textbf{Shell Companies:} Transfer funds to shell company with no operations
				\item \textbf{Private Loans:} Loan to shell company matching criminal proceeds
				\item \textbf{Multiple Transfers:} Move funds through multiple accounts
				\item \textbf{Asset Purchases:} Buy assets (art, real estate, crypto)
				\item \textbf{False Invoices:} Use false invoices to justify transfers
				\item \textbf{Complex Structures:} Use trusts, offshore companies
			\end{itemize}
		\end{block}
		\vspace{0.02cm}
		\begin{block}{\fontsize{7}{9}\selectfont \textbf{Identifying Layering — CAMS Exam Favorite}}
			\scriptsize
			\textbf{Scenario:} Private loan extended to shell company with no operations. Loan amount matches criminal proceeds.
			\par \textbf{Classification:} \textbf{Layering tactic} — designed to conceal origin of illicit funds.
		\end{block}
		\vspace{0.02cm}
		\recallbox{\scriptsize \textbf{Key Rule:} Layering = \textbf{concealing origin}. Placement = \textbf{entering system}. Integration = \textbf{using funds legitimately}.}
	\end{frame}
	
	% ================================================================
	% SLIDE 31: FINAL EXAM TIPS - DOMAIN A
	% ================================================================
	\begin{frame}
		\frametitle{\fontsize{11}{13}\selectfont \textbf{Domain A — Critical Exam Tips}}
		\scriptsize
		\begin{block}{\fontsize{7}{9}\selectfont \textbf{Most Common Traps}}
			\scriptsize
			\begin{itemize}
				\item \textbf{Rec 15 vs Rec 16:} Rec 15 = VASP framework. Rec 16 = Travel Rule.
				\item \textbf{Black vs Grey List:} Black = countermeasures. Grey = EDD.
				\item \textbf{Decentralized vs Centralized:} Decentralized = anonymity. Centralized = KYC.
				\item \textbf{SOW vs SOF:} SOW = overall wealth. SOF = specific funds.
				\item \textbf{De-Risking:} \textbf{Not} a substitute for managing risk. FATF discourages blanket de-risking.
				\item \textbf{Tipping Off:} \textbf{Never} disclose SAR/STR investigation to customer.
				\item \textbf{Structuring:} Illegal even with clean money.
				\item \textbf{SDD:} Suspicion overrides — must do full CDD if suspicious.
				\item \textbf{Shell Banks:} Always prohibited.
				\item \textbf{Rec 17:} Can rely on third parties, but \textbf{remain liable}.
			\end{itemize}
		\end{block}
		\vspace{0.02cm}
		\begin{block}{\fontsize{7}{9}\selectfont \textbf{Exam Strategy}}
			\scriptsize
			\begin{itemize}
				\item \textbf{Read Carefully:} "Select Two" vs "Select Three" vs "Select All"
				\item \textbf{Eliminate Wrong Answers:} Look for absolute statements (always, never)
				\item \textbf{Think Risk-Based:} When in doubt, apply the risk-based approach
				\item \textbf{Document:} Always document your decisions and rationale
				\item \textbf{Escalate:} When in doubt, escalate to compliance
			\end{itemize}
		\end{block}
		\vspace{0.02cm}
		\recallbox{\scriptsize \textbf{Final Rule:} Compliance is about \textbf{understanding the risk} and applying \textbf{proportionate controls}.}
	\end{frame}
	
	% ================================================================
	% END SLIDE
	% ================================================================
	\begin{frame}
		\centering
		\vspace{1.5cm}
		{\fontsize{18}{22}\selectfont \textbf{Domain A — Complete}}\\[0.2cm]
		{\fontsize{11}{13}\selectfont All Hard Concepts with Detailed Explanations}\\[0.1cm]
		\rule{4cm}{0.5pt}\\[0.1cm]
		{\fontsize{8}{10}\selectfont VASP · Crypto · FATF Recommendations · Black/Grey Lists}\\[0.05cm]
		{\fontsize{8}{10}\selectfont Investment Banking · Wealth Management · PSP · OSINT · KYC}\\[0.05cm]
		{\fontsize{8}{10}\selectfont TCSP · Data Governance · Post-Acquisition · Sanctions · DORA}\\[0.05cm]
		{\fontsize{8}{10}\selectfont Transaction Monitoring · Gambling · Asset Management · Fraud}\\[0.05cm]
		{\fontsize{8}{10}\selectfont TBML · Structuring · Layering · FTZ · Third-Party Reliance}\\[0.1cm]
		\rule{4cm}{0.5pt}\\[0.1cm]
		{\fontsize{8}{10}\selectfont \textit{All traps explained. All concepts covered. Good luck on your CAMS exam!}}
		\vspace{0.5cm}
	\end{frame}
	
\end{document}